% ============================================================
% CAPITOLUL 7 — APLICAȚIA WEB E-MEDICINE
% ============================================================

\chapter{Aplicația Web E-Medicine}
\label{ch:aplicatie}

\section{Rolul aplicației}
\label{sec:functionalitati}

Aplicația web este interfața prin care modelul matematic și infrastructura de
execuție devin utilizabile. Ea nu este un sistem clinic complet, ci un mediu de
lucru pentru încărcarea datelor, configurarea simulărilor SIDARTH sau SIR și
compararea rezultatelor obținute în R, Julia și Python. Din această perspectivă,
aplicația are rol de instrument experimental.

Fluxul principal începe cu autentificarea utilizatorului. După autentificare,
acesta poate încărca un fișier CSV cu date epidemiologice, poate alege modelul,
poate introduce parametrii și poate urmări statusul execuției. Când runner-ul
finalizează job-ul, interfața afișează traiectoriile compartimentale și permite
exportul rezultatelor în format CSV. O pagină separată permite compararea mai
multor rulări finalizate, iar pagina de benchmark prezintă agregările calculate la
nivel de limbaj.

\section{Backend Go}
\label{sec:backend_go}

Backend-ul este implementat în Go și expune un API REST sub prefixul
\texttt{/api}. Serverul folosește routerul \texttt{chi}, middleware pentru
logging, recuperare după panică, timeout de request și CORS, iar accesul la rutele
protejate este controlat prin JWT. Alegerea Go este potrivită pentru această
componentă deoarece API-ul este stateless, are nevoie de latențe predictibile și
poate fi distribuit ușor sub formă de binar static în container.

Structura codului separă intrarea HTTP de accesul la date. Pachetul
\texttt{handlers} parsează cererile, validează input-ul și construiește
răspunsurile JSON. Pachetul \texttt{repository} conține interogările SQL și
operațiile tranzacționale peste \texttt{pgx}. Pachetul \texttt{runner} definește
schema mesajului trimis în Redis și publicarea job-urilor. Contractul de simulare
include tipul modelului, parametrii, condițiile inițiale și conexiunea la baza de
date.

\section{Gestionarea dataseturilor}
\label{sec:dataseturi_web}

Dataseturile sunt încărcate printr-un formular multipart. API-ul acceptă câmpurile
\texttt{name}, \texttt{disease\_code} și \texttt{file}; fișierul trebuie să conțină
cel puțin coloanele \texttt{date}, \texttt{region}, \texttt{infected} și
\texttt{population}. Coloanele \texttt{recovered} și \texttt{susceptible} sunt
opționale, dar atunci când sunt prezente sunt folosite pentru condițiile inițiale.
Pentru SIDARTH, cazurile active sunt plasate inițial în compartimentul $I$, iar
persoanele ieșite din infecție în $H$.

Persistarea datasetului se face într-o tranzacție. Mai întâi este creat rândul din
\texttt{datasets}, apoi înregistrările epidemiologice sunt inserate prin
\texttt{CopyFrom}. Această alegere reduce costul încărcării CSV-urilor față de
inserarea rând-cu-rând și păstrează atomicitatea: dacă una dintre înregistrări nu
poate fi scrisă, datasetul nu rămâne într-o stare parțială.

\section{Execuția simulărilor}
\label{sec:executie_simulari_web}

Crearea unei simulări este un proces scurt în API și un proces asincron în
ansamblul sistemului. Endpointul \texttt{POST /api/simulations} verifică existența
datasetului, modelul, limbajul selectat, parametrii și $t_{\max}$, apoi extrage
prima înregistrare epidemiologică pentru a determina valorile inițiale. După ce
rândul din \texttt{simulation\_runs} este creat cu status \texttt{pending}, API-ul
publică mesajul JSON în coada Redis corespunzătoare limbajului.

Runner-ele sunt responsabile pentru tranzițiile ulterioare. Ele marchează rularea
ca \texttt{running}, execută solver-ul pentru \texttt{model\_type}, inserează
rezultatele și benchmarkul, apoi marchează rularea ca \texttt{completed}. Dacă
apare o eroare, statusul devine \texttt{failed}, iar mesajul este salvat în baza de
date. Interfața web nu trebuie să țină o conexiune deschisă pe durata calculului;
pagina de rezultat interoghează periodic endpointul de status până la finalizare.

\section{Frontend Vue.js}
\label{sec:frontend_vue}

Frontend-ul este o aplicație Vue 3 construită cu Vite și TypeScript. Rutarea este
gestionată cu Vue Router, iar starea aplicației cu Pinia. Clientul HTTP este o
instanță Axios configurată cu prefixul \texttt{/api}; un interceptor adaugă automat
headerul \texttt{Authorization} atunci când există token în \texttt{localStorage}.

Paginile sunt organizate după fluxurile principale ale aplicației:
\texttt{DashboardPage} sumarizează dataseturile și simulările recente,
\texttt{DatasetsPage} gestionează încărcarea CSV-urilor,
\texttt{NewSimulationPage} colectează modelul și parametrii,
\texttt{SimulationResultPage} afișează statusul și rezultatele,
\texttt{ComparePage} suprapune seriile mai multor rulări, iar
\texttt{BenchmarksPage} afișează agregările pe model și limbaj. Componenta
\texttt{CompartmentChart} folosește Chart.js prin \texttt{vue-chartjs} și
construiește dinamic seriile în funcție de cheile disponibile în
\texttt{compartments}. Astfel, același grafic poate afișa atât SIDARTH, cât și SIR.

Decizia de a folosi polling la două secunde este simplă, dar adecvată pentru
durata așteptată a simulărilor. O soluție bazată pe WebSocket sau Server-Sent
Events ar reduce numărul de cereri, însă ar introduce stare suplimentară în
backend. Pentru un API stateless și pentru job-uri care durează de la câteva
secunde la câteva minute, polling-ul păstrează implementarea previzibilă.

\section{Validare funcțională}
\label{sec:testare_functionala}

Validarea funcțională a aplicației poate fi făcută prin fluxuri end-to-end care
urmează operațiile reale ale utilizatorului: înregistrare, autentificare, încărcare
dataset, pornire simulare SIDARTH, așteptare finalizare, afișare grafic și export
CSV. Aceeași suită ar trebui să verifice și rularea SIR, deoarece aceasta asigură
compatibilitatea cu modelul secundar și cu comparațiile simple.
